\documentclass[10pt,a4paper,titlepage,twocolumn]{scrartcl} % Professionelle KOMA-Klasse
% \usepackage{tgpagella}

% \usepackage{tgbonum}
% \usepackage{XCharter} % Setzt XCharter als Hauptschrift
% \usepackage[utf8]{inputenc}
\usepackage[T1]{fontenc}
\usepackage{ebgaramond}
\usepackage[english]{babel}
\usepackage{longtable} 
\usepackage{booktabs}   
\usepackage[table]{xcolor}
\usepackage{geometry}
\usepackage{fontawesome5}
\usepackage{tcolorbox} % Für minimalistische, schicke RPG-Rahmen und Boxen
\usepackage{fontspec}
\usepackage{imakeidx}


\makeindex[columns=2, title=Index, intoc] 

% \setmainfont{TeX Gyre Pagella}
\newfontfamily\dsix{DpolySixSider-jqRO.ttf}

% Seitenränder perfekt für ein ausgedrucktes RPG-Heft einstellen
\geometry{top=2.5cm, bottom=2.5cm, left=2.5cm, right=2.5cm}

% Eine eigene, schlichte Box für RPG-Hinweise definieren (Schwarz/Weiß-Stil)
\newtcolorbox{rpgbox}[1]{
  colback=white, 
  colframe=black, 
  fonttitle=\bfseries, 
  title=#1, 
  arc=0mm, % Eckige Kanten für Retro-Look
  boxrule=1pt
}


\newcommand{\kw}[1]{\textsf{\textbf{#1}}\index{#1}}
\newcommand{\game}{\textsf{\textbf{Beneath The Shard}}}
\newcommand{\todo}{\textit{todo}}
\newcommand{\world}{\textit{WORLDNAME}}


% ==========================================
% TITELBLATT CONFIGURATION (Manuell & Unabhängig)
% ==========================================5
\begin{document}

\begin{titlepage}
    \centering
    \vspace*{3cm}
    {\Large\bfseries A Solo Pen \& Paper RPG  \par}
    \vspace{1cm}
    {\Huge\bfseries Beneath The Shard \par}
    \vspace{0.5cm}
    {\large\bfseries -- Chronicles of the Star-Grave -- \par}
    \vspace{2cm}
    
    \begin{minipage}{0.7\textwidth}
        \centering
        \itshape
        A cosmic catastrophe shattered a distant world of magic. Now, its glowing 
        remnants grow deep beneath our feet, carving twisting dungeons and warping 
        reality itself. Steel your heart, adventurer -- the Shard calls.
    \end{minipage}
    
    \vfill
    {\large Planet Terramyr Campaign Setting \par}
    \vspace{0.5cm}
    {\large Version 1.0 -- \today \par}
\end{titlepage}


\tableofcontents

\clearpage 

\section{Playing The Game}

To start a new game of \game{} you need:

\begin{itemize}
    \item \textbf{Dice:} at least 12 six-sided dice, and two ten-sided dice of different color.
    \item \textbf{Party Sheet:} \todo{}
    \item \textbf{World Map:} \todo{}
    \item \textbf{Scrap Paper:} \todo{}
\end{itemize}

First you create a party, then you get on a journey through \world{}. On your your journey you will
face combat- and non-combat encounters.

\section{Character Creation}

In \game{} you play a party of three characters. Each character is created based on an \kw{archetype}.
The three archetypes are:

\begin{itemize}
    \item \textbf{Warrior:} \todo{}
    \item \textbf{Rogue:} \todo{}
    \item \textbf{Mage:} \todo{}
\end{itemize}

\subsection{Attributes}

Each character has six attributes:
\begin{itemize}
    \item \textbf{Strength     $\left[\textbf{STR}\right]$:} \todo{}
    \item \textbf{Dexterity    $\left[\textbf{DEX}\right]$:} \todo{}
    \item \textbf{Constitution $\left[\textbf{CON}\right]$:} \todo{}
    \item \textbf{Intelligence $\left[\textbf{INT}\right]$:} \todo{}
    \item \textbf{Wisdom       $\left[\textbf{WIS}\right]$:} \todo{}
    \item \textbf{Charisma     $\left[\textbf{CHA}\right]$:} \todo{}
\end{itemize}

\subsection{Attribute Value}
Each attribute has an \kw{Attribute Value} ranging from \textsf{0} to \textsf{100}. 

Roll \textsf{6} times for each 
of your character's attributes on the following table with \D{3}{6}

\begin{table}[htbp]
  \centering % Centers the entire table on the page
  \caption{Determine Attribute Values}
  \label{tab:attribute_values} % MUST come after or inside \caption
\renewcommand{\arraystretch}{1.3}
\begin{tabular}{|c|c|}
    \hline
    \textbf{3d6} & \textbf{Attribute Value} \\
    \hline
    \textsf{2 -- 4}   &  \textsf{20} \\
    \hline                    
    \textsf{5 -- 6}   &  \textsf{30} \\
    \hline                
    \textsf{7 -- 8}   &  \textsf{40} \\
    \hline             
    \textsf{9 -- 10}  &  \textsf{50} \\
    \hline         
    \textsf{11 -- 12} &  \textsf{60} \\
    \hline        
    \textsf{13 -- 14} &  \textsf{70} \\
    \hline    
    \textsf{15 -- 16} &  \textsf{80} \\
    \hline
    \textsf{17 -- 18} &  \textsf{90} \\
    \hline
\end{tabular}
\end{table}

\subsection{Attribute Rank}

The \kw{Attribute Rank} of an attribute is its value divided by \textsf{10}.
If a character's \attr{STR} is \textsf{75}, its rank is \textsf{7}

\subsection{Attribute Check}

Sometimes the game asks you to perform an attribute check. Then you need to roll 
d100 equal or below your attribute value.

\begin{rpgboxb}{Attribute Check Example}
    \textit{Check \acheck{DEX}{+10}:} you need add \textsf{10} to your \attr{DEX} attribute
    and roll d100 equal or below it.

    If a character has \attr{DEX} \textsf{60} he needs to roll \textsf{70} or less on a \dper
\end{rpgboxb}






\subsection{Party}


\clearpage


\printindex

\end{document}